\section{Centrality and rational stability}\label{sec:stability}

Throughout this section, let $R=\F_2\oplus I$ be a unital associative
algebra, where $I$ is a two-sided ideal and the indicated copy of
$\F_2$ contains the identity of $R$. Suppose that every matrix over
$I$ is nilpotent. These are among the
properties supplied by Theorem~\ref{thm:algebra}. No commutativity or finite
generation assumption on $R$ is needed in this Section. Write
$H_n=\GL_n(R)$, retain the groups $\E_n(R)$ and $\St_n(R)$ of
Section~\ref{sec:finite-presentation}, and set
\[
 J_n=\Ker\bigl(\St_n(R)\longrightarrow\E_n(R)\bigr).
\]
All maps between successive ranks are the standard coordinate
stabilization maps.
Our two outputs are centrality of $J_n$ and extension of every rational
homomorphism on $J_n$ to larger ranks. The latter will allow torsion
proved after stabilization to descend to rank twelve.

\subsection{Frames and presentations of the full matrix groups}

\begin{lemma}\label{st:linear-algebra}
A square matrix over $R$ is invertible if and only if its reduction modulo
$I$ is invertible. Every tuple of columns whose reductions are linearly
independent extends to an invertible matrix. The group $H_n$ is generated
by elementary matrices and the diagonal matrices $d_i(u)$ having a unit
$u\in R^\times$ in one position and ones elsewhere.
\end{lemma}

\begin{proof}
The invertibility assertion is Lemma~\ref{fp:matrix-periodicity}.
To complete a tuple of columns, complete their reductions to a basis
over $\F_2$ and lift the additional columns; the resulting matrix is
invertible by that assertion.

In particular, an element of $R$ with nonzero reduction is a unit. An
invertible matrix has such an entry in its first column. Interchange
rows to put that entry in position $(1,1)$, and denote it by $u$.
Left multiplication by elementary matrices with parameters
$-a_{i1}u^{-1}$ clears the remaining entries of the first column.
Right multiplication by elementary matrices with parameters
$-u^{-1}a_{1j}$ then clears the remaining entries of the first row.
The resulting complementary block is invertible by reduction. Induction
reduces the matrix to a diagonal matrix of units. The order of the
products used here is valid in a noncommutative ring. Finally, a row
interchange is elementary because, in characteristic two, its two by two
block is
\[
 e_{12}(1)e_{21}(1)e_{12}(1)=
 \begin{pmatrix}0&1\\1&0\end{pmatrix}.
\]
This proves the generation assertion.
\end{proof}

A \emph{frame} is a tuple of columns of $R^n$ with linearly independent
reductions in $\F_2^n$. Let $X_n$ be the abstract simplicial complex whose
vertices are columns with nonzero reduction and whose simplices are the
finite sets that form frames. Distinct columns remain distinct vertices
even when their reductions coincide.

Complexes of partial bases and frames are classical tools in
homological stability. Van der Kallen treats associative rings under
Bass stable-range conditions \cite{vanderKallen1980Homology}. We prove
the connectivity and rational comparison needed for our particular
ring directly.

\begin{lemma}\label{st:frames}
The complex $X_n$ has vanishing reduced rational homology in degrees at
most $n-2$. It is simply connected for $n\geq3$.
\end{lemma}

\begin{proof}
We first work with a finite labeled collection of nonzero vectors
spanning a vector space of dimension $n$ over $\F_2$. Its independence
complex has one vertex for every label; parallel vectors with different
labels are not identified. We prove the homology assertion by induction
on $n$. For $n=1$, the asserted vanishing is just nonemptiness.

Choose a basis from the collection. Its independence complex is a
simplex. Adjoin the remaining labeled vectors one at a time. When a
vector $v$ is adjoined, its link in the preceding complex is the
independence complex of the preceding vectors projected to the quotient
by $\F_2v$, with zero projections discarded. The projected collection
spans a space of dimension $n-1$, and its labels are retained.
By induction its reduced homology vanishes through degree $n-3$.
Adjoining $v$ attaches the cone on this link. The relative simplicial
chains are its augmented link chains shifted by one. The long exact
sequence of this pair therefore preserves vanishing of reduced homology
through degree $n-2$.

For $n\geq3$ each of these links is connected, by the homology assertion
in rank $n-1$. Attaching a cone on a connected subcomplex to a simply
connected complex preserves simple connectivity: a path through the
new cone may be replaced, with the same endpoints, by a path in its
link, and the replacement is homotopic relative to the endpoints in
the cone. Starting with the basis simplex proves the finite assertion
on fundamental groups as well.

Any finite set of vertices of $X_n$ can be augmented by columns whose
reductions form a basis. The independence complex on that finite set
is one of the finite complexes just considered. Every simplicial cycle
and every edge loop uses finitely many vertices, so the finite results
give the corresponding assertions for $X_n$.
\end{proof}

To prove centrality, we will assign to each elementary matrix
conjugation by its Steinberg root and extend this assignment to an
action of $H_n$ on $\St_n(R)$. For $n\geq6$, a matrix relation on at
most five coordinates leaves a spare coordinate on which to check the
action. The next proposition reduces all matrix relations to relations
of that form.

Call a word in elementary and single-position diagonal generators
\emph{supported on $S$} if every index occurring in it belongs to $S$.

\begin{proposition}\label{st:small-presentation}
For $n\geq5$, the group $H_n$ has the presentation with elementary
generators $e_{ij}(a)$ and diagonal generators $d_i(u)$, subject to all
true matrix relations supported on coordinate sets of size at most
five.
\end{proposition}

\begin{proof}
Let $H'_n$ be the presented group and let $\pi:H'_n\to H_n$ be its
natural surjection, which is surjective by
Lemma~\ref{st:linear-algebra}. We prove injectivity by induction on $n$.
For $n=5$, every true relation is included, so there is nothing to prove.
Suppose that $n>5$ and the result is known in rank $n-1$.

The stabilizer $P$ of the column $e_1$ in $H_n$ consists of matrices
\[
 \begin{pmatrix}1&b\\0&B\end{pmatrix},
 \qquad B\in H_{n-1}.
\]
Its complementary copy of $H_{n-1}$ lifts homomorphically to $H'_n$ by
the induction hypothesis. Its row subgroup, generated by $e_{1j}(a)$
for $j\ne1$, also lifts: its additive and commuting relations involve
at most three coordinates. Conjugation of this row subgroup by each
complementary elementary or diagonal generator is its matrix
conjugation, again by relations on at most three coordinates. These
relations give a homomorphic section $P\to H'_n$ of $\pi$ on $P$.
For example, writing a matrix of $P$ uniquely as a row transvection
times $\operatorname{diag}(1,B)$ gives its lift using this semidirect
product structure. Denote the image by $P'$. The section is injective
because its composite with $\pi$ is the inclusion of $P$.

Let $s\in H'_n$ be the elementary word for the swap $(12)$. Its square
is one by a two-coordinate relation. The groups $P'$ and $\langle
s\rangle$ generate $H'_n$. Indeed, conjugating the generators already
in $P'$ by $s$ gives every first-column elementary generator and the
first diagonal generator; these conjugation identities involve at
most three coordinates.

Let $Q\leq P$ fix both $e_1$ and $e_2$ pointwise. It is generated by
the complementary elementary and diagonal generators on coordinates
$3,\ldots,n$ and the row generators $e_{ij}(a)$ with $i\in\{1,2\}$
and $j\geq3$. Its image $Q'$ under the section of $P$ is generated by
the same displayed generators in $H'_n$. Conjugation by $s$ fixes the
complementary generators and exchanges the two row families, by
small-coordinate relations. Hence
\[
 sQ's^{-1}=Q'.
\]
At the first induction step $n=6$, this argument uses only generators
of the complementary $H_4$ inside the already lifted $P$; it does not
require a presentation theorem in rank four.

We use $X_n$ to identify the cosets of $P'$. The map
\[
 H'_n/P'\longrightarrow X_n^{(0)},\qquad
 hP'\longmapsto\pi(h)e_1
\]
is well defined. If $w$ is adjacent to $\pi(h)e_1$, frame completion
provides $p\in P$ with $p e_2=\pi(h)^{-1}w$. From $hP'$ transport
along this oriented edge to
\[
 h\widetilde p\,sP',
\]
where $\widetilde p$ is the section of $p$. The result is independent
of $p$: if $p e_2=p'e_2$, then $p'^{-1}p\in Q$, and
$s^{-1}Q's\subseteq P'$. It is also independent of the representative
$h$ of its coset, since replacing $h$ by $hr$, with $r\in P'$, may
be offset by replacing $p$ by $\pi(r)^{-1}p$. None of these assertions
assumes global injectivity of $\pi$.

Transport reverses along an edge: from the representative
$h\widetilde p s$, the reverse step can use $p=1$, giving
$h\widetilde p s^2P'=hP'$. Transport around every triangle is also
trivial. By translating its first vertex and then using one element
of $P$, frame completion reduces the triangle to $(e_1,e_2,e_3)$.
Put $t=(23)\in P'$, expressed as an elementary word. The three
successive representatives can be taken to be
\[
 s,\qquad sts,\qquad ststs=t.
\]
The final equality is a true relation on the first three coordinates,
and $t\in P'$. Thus the endpoint coset is the starting coset.

By Lemma~\ref{st:frames}, $X_n$ is simply connected. Combinatorial
edge paths with the same endpoints therefore differ by backtracks
and triangle moves, after subdivisions of a null homotopy; both moves
preserve transport. Starting at $P'$, transport consequently assigns
exactly one coset to every vertex. Every coset is reached, since
$P'$ and $s$ generate $H'_n$: multiplication on the right by an
element of $P'$ changes only a representative, and multiplication by
$s$ is one of the allowed edge steps. The displayed coset map is
therefore a bijection.

If $h\in\Ker\pi$, its coset maps to $e_1$, so $h\in P'$.
The restriction of $\pi$ to $P'$ is injective. Hence $h=1$, completing
the induction.
\end{proof}

\subsection{The central kernel and the quotient by elementary matrices}

The faithful-strip argument below is a finite-rank form of the
argument in Steinberg's proof of stable centrality
\cite[Section~7, Theorem~14]{Steinberg1967}. The presentation on small
coordinate sets provides the rank control here. Centrality for the
present ring also lies within Voronetsky's semilocal-ring theorem
\cite[Theorem~1]{Voronetsky2021Centrality}: the nonunits of $R$ are
exactly $I$, since each $1+a$ with $a\in I$ is invertible and a nilpotent
element cannot be a unit. Thus $R$ is local, and $M_n(R)$ is semilocal
with the standard matrix idempotents. We retain the direct strip proof
and its stated rank bound.

\begin{proposition}\label{prop:centrality}
For $n\geq6$, the subgroup $J_n$ is central in $\St_n(R)$.
\end{proposition}

\begin{proof}
Assign to each elementary generator $e_{ij}(a)$ of $H_n$ conjugation
by $x_{ij}(a)$ on $\St_n(R)$. Assign to a diagonal generator the
automorphism
\[
 x_{ij}(a)\longmapsto x_{ij}(u_i a u_j^{-1}),
\]
where all but one of the $u_i$ are one. This preserves the Steinberg
relations; in particular the adjacent-root product becomes
$(u_i a u_j^{-1})(u_j b u_k^{-1})=u_i ab u_k^{-1}$.

We check that every true relation supported on a set $S$ of at most
five coordinates acts trivially. Choose $t\notin S$. For parameters
indexed by $S$, the two groups
\[
 U_t(v)=\prod_{i\in S}x_{it}(v_i),\qquad
 V_t(w)=\prod_{i\in S}x_{ti}(w_i)
\]
are additive abelian groups. Their maps to elementary matrices are
injective, since their parameters occur in distinct matrix positions
and no cross product occurs inside either strip.

The assigned action of $e_{ab}(r)$, with $a,b\in S$, sends
$v_a$ to $v_a+r v_b$ in the first strip. In the second it sends
$w_b$ to $w_b-w_a r$. To see the latter formula directly, use
\[
 [x_{ta}(w_a),x_{ab}(r)]=x_{tb}(w_a r);
\]
the right side commutes with both factors, so reversing the
commutator replaces its parameter by its negative. A diagonal
$d_a(u)$ sends $v_a$ to $u v_a$ and $w_a$ to $w_a u^{-1}$.
Thus a word with matrix product $M$ supported on $S$ acts on these
two strips by
\[
 v\longmapsto Mv,\qquad w\longmapsto wM^{-1}.
\]
These formulas preserve the order of all noncommutative products.
A true relation has $M=\Id$, so fixes both strips pointwise.
Repeating this calculation for each $t\notin S$ fixes every root
with exactly one index in $S$. It also fixes every root supported
outside $S$. Finally, for $i,j\in S$,
\[
 x_{ij}(a)=[x_{it}(a),x_{tj}(1)],
\]
so it fixes the internal roots as well. Exactly one outside coordinate
is needed for this argument.

Proposition~\ref{st:small-presentation} now gives a homomorphism
$H_n\to\operatorname{Aut}(\St_n(R))$. On the image of each
Steinberg generator it is conjugation by that generator. Hence for
every $z\in\St_n(R)$, the action of its matrix image is conjugation
by $z$. If $z\in J_n$, that matrix image is the identity. Therefore
$z$ commutes with every element of $\St_n(R)$.
\end{proof}

\begin{lemma}\label{st:quotient}
The group $\E_n(R)$ is normal in $H_n$. For $n\geq5$, inclusion
induces an isomorphism
\[
 H_n/\E_n(R)\longrightarrow H_{n+1}/\E_{n+1}(R).
\]
\end{lemma}

\begin{proof}
Elementary generators normalize their own generated subgroup, and
diagonal conjugation preserves each elementary root subgroup.
Lemma~\ref{st:linear-algebra} therefore proves normality.

In the presentation of Proposition~\ref{st:small-presentation}, kill
all elementary generators. Elementary swap words then become trivial,
so the conjugation relation for a swap identifies $d_i(u)$ with
$d_j(u)$ for every pair of positions. Diagonals in different positions
commute. We may consequently use a single generator $d(u)$ for each
unit. Every remaining defining relation is obtained from a true
matrix relation on at most five coordinates by deleting elementary
generators and replacing each $d_i(u)$ by $d(u)$.

This list of relations is independent of $n\geq5$. Indeed, whether a
word supported on $S$ is a true relation depends only on its matrix
block on $S$, and not on the other coordinates. Every such word can
be relabeled onto at most five fixed coordinates, and all those
relations already occur in rank five. Inclusion is the identity on
the resulting unit generators and on their relations, proving the
assertion.
\end{proof}

\subsection{Central extensions and the rational dual of the kernel}

We now turn to extension of homomorphisms $J_n\to\Q$. We first
identify these homomorphisms with central extensions of $\E_n(R)$
by $\Q$, and then with the dual of its second rational homology.
This comparison will reduce the desired extension property to a
homology calculation for the frame complex. The connection between
central extensions, second homology, and Steinberg kernels is classical;
see \cite[Section~7]{Steinberg1967} and
\cite[Sections~III.5.3--III.5.5]{WeibelKbook}. Here the comparison
must also preserve the specified maps between ranks.

We use rational group homology with arbitrary left coefficient modules.
It may be computed as
\[
 H_q(G;M)=H_q\bigl(P_\bullet\otimes_{\Q[G]}M\bigr),
\]
where $P_\bullet\to\Q$ is a free right $\Q[G]$-resolution. For
example, the homogeneous bar resolution has ordered lists of group
elements as generators, deletion boundary, and free diagonal group
action. Its augmented complex is exact by the contraction that inserts
the identity element. The usual comparison maps between free
resolutions are obtained by lifting successive generators through
surjective boundary maps; the same procedure gives the comparison
homotopies.

\begin{lemma}\label{st:finite-splitting}
Every central extension of a locally finite group by the additive
group $\Q$ has a unique homomorphic splitting.
\end{lemma}

\begin{proof}
First let the quotient group $L$ be finite. Choose normalized
set-theoretic lifts and write their additive multiplication defect
as $f(g,h)\in\Q$. Associativity says
\[
 f(g,h)+f(gh,t)=f(h,t)+f(g,ht).
\]
Define $b(g)=|L|^{-1}\sum_{t\in L}f(g,t)$. Averaging this identity
over $t$ gives
\[
 f(g,h)=b(g)+b(h)-b(gh).
\]
Changing the chosen lift of $g$ by the central element $-b(g)$
therefore makes it a homomorphic splitting. Two splittings differ
by a homomorphism $L\to\Q$, which must be zero.

If $L$ is locally finite, apply the finite result to the inverse
image of each finite subgroup. Uniqueness makes the splittings
compatible under inclusion. Every finite set of elements of $L$
lies in a finite subgroup, so these splittings glue to a unique
homomorphic splitting on $L$.
\end{proof}

\begin{lemma}\label{st:extensions}
For $n\geq6$, equivalence classes of central extensions of
$\E_n(R)$ by the additive group $\Q$ are naturally in bijection
with $\Hom(J_n,\Q)$. This bijection commutes with restriction
along inclusions of ranks.
\end{lemma}

\begin{proof}
Every upper unitriangular subgroup over $R$, for any ordering of
the coordinates, is locally finite. To see this, start with
finitely many unitriangular matrices and collect their entries
above the diagonal in a finite set $S$. Entries of products and
inverses are prime-field linear combinations of words in $S$
of length at most $n-1$. Indeed, every such product follows a
strictly increasing sequence of coordinate indices; the inverse
formula is the finite geometric series for a strictly triangular
matrix and has the same property. There are finitely many words
of these bounded lengths, and their $\F_2$-span is finite.
Only finitely many matrices can consequently occur in the
subgroup generated by the original matrices.

Let
\[
 1\longrightarrow\Q\xrightarrow{\iota}X
 \longrightarrow\E_n(R)\longrightarrow1
\]
be a central extension. Lemma~\ref{st:finite-splitting} gives a
unique splitting on each unitriangular subgroup. Their root
lifts agree whenever a root occurs in different orderings: the
root group itself is locally finite and has a unique splitting.
These root lifts satisfy every Steinberg relation. Addition is
a relation in one root group. A consecutive-root commutator
relation fits in a unitriangular group with its three indices
in the path order. A commuting relation fits in a unitriangular
group as well: the two directed edges specified by its root
positions have no directed cycle, since the opposite-root case
is excluded. Extend the resulting partial ordering to an
ordering of all coordinates. In each case the relevant
splitting is a homomorphism, so proves the relation among lifts.
We have obtained a homomorphism
\[
 f:\St_n(R)\longrightarrow X
\]
over $\E_n(R)$. Its restriction to $J_n$ lands in the prescribed
central copy of $\Q$ and defines $\varphi:J_n\to\Q$.
The map $f$ is unique: two such maps differ pointwise by a
homomorphism from the perfect group $\St_n(R)$ to the central
abelian group $\Q$. Perfectness also follows directly from
$x_{ij}(a)=[x_{ik}(a),x_{kj}(1)]$.

Conversely, given $\varphi\in\Hom(J_n,\Q)$, form
\[
 X_\varphi=
 (\St_n(R)\times\Q)/
 \{(j,-\varphi(j)):j\in J_n\}.
\]
The subgroup being divided out is central by
Proposition~\ref{prop:centrality}. The map $q\mapsto[(1,q)]$
is injective, and the induced map $X_\varphi\to\E_n(R)$
is surjective with this central kernel. In fact
$[(j,q)]=[(1,q+\varphi(j))]$ for $j\in J_n$.
Thus this is a central extension by the specified copy of
$\Q$. The map $s\mapsto[(s,0)]$ gives its canonical root
lifts, and its restriction to $J_n$ recovers $\varphi$.

Starting instead with the extension $X$ and its map $f$, the
homomorphism
\[
 \St_n(R)\times\Q\longrightarrow X,\qquad
 (s,q)\longmapsto f(s)\iota(q)
\]
is surjective: choose $s$ with the same elementary image as
any given element of $X$, and account for the difference in
the central kernel. Its kernel is exactly
$\{(j,-\varphi(j)):j\in J_n\}$. It therefore identifies
$X_\varphi$ with $X$ as an extension. This proves the
bijection. Restricting an extension to a smaller rank
preserves its canonical root lifts, by their uniqueness.
The bijection consequently commutes with rank restriction.
\end{proof}

For completeness, the relation between these extensions and
rational homology can also be read directly from chains. For any
group $G$, a choice of normalized lifts in a central extension
by $\Q$ gives a function $f:G\times G\to\Q$ satisfying
\[
 f(g,h)+f(gh,k)=f(h,k)+f(g,hk).
\]
Changing the lifts adds a coboundary
$b(g)+b(h)-b(gh)$. Conversely, a normalized function satisfying
this identity defines an extension on $\Q\times G$ with
product
\[
 (a,g)(b,h)=(a+b+f(g,h),gh).
\]
The identity gives associativity, and also the required inverse
identities. Thus extension classes are the degree-two
cohomology of the rational bar cochain complex. This is the
linear dual of the bar chain complex computing $H_*(G;\Q)$.
Every exact sequence of vector spaces splits, so evaluation
on cycles gives a natural isomorphism
\[
 H^2(G;\Q)\cong
 \Hom_{\Q}\bigl(H_2(G;\Q),\Q\bigr).
\]
The isomorphism is natural even though a splitting can be used
to prove that it is bijective. Combining this identification with
Lemma~\ref{st:extensions} gives, for $n\geq6$, the natural bijections
\[
 \Hom(J_n,\Q)\longleftrightarrow H^2(\E_n(R);\Q)
 \cong\Hom_{\Q}\bigl(H_2(\E_n(R);\Q),\Q\bigr).
\]
Thus it remains to show that the standard inclusion induces an
isomorphism on second rational homology for $n\geq12$. Consequently
every homomorphism $J_n\to\Q$ will extend to $J_{n+1}$.

\subsection{Rational homology of stabilizers}

We prove the required homology comparison using the action on the frame
complex $X_m$.
The stabilizers of unordered frames contain permutations of the frame.
For frames of two, three, or four columns, the orientation signs will
force the stabilizer homology to vanish, leaving only the vertex
stabilizer in the degree-two calculation.

\begin{lemma}\label{st:coinvariants}
Let $N$ be a locally finite normal subgroup of $K$, and let $M$ be a
left $\Q[K]$-module. There are natural isomorphisms
\[
 H_q(K;M)\cong H_q(K/N;M_N),
\]
where $M_N$ denotes the module of $N$-coinvariants.
\end{lemma}

\begin{proof}
First, taking $N$-coinvariants is exact over $\Q$. Only preservation
of injections needs proof. If $U\subseteq V$ and $u\in U$ becomes
zero in $V_N$, write
\[
 u=\sum_i(n_i v_i-v_i).
\]
The subgroup $B$ generated by the finitely many $n_i$ is finite.
The averaging operator $e_B=|B|^{-1}\sum_{b\in B}b$ kills the
right side, so $e_Bu=0$. Hence
\[
 u=u-e_Bu=|B|^{-1}\sum_{b\in B}(u-bu)
\]
is already zero in $U_N$.

Take a free left $\Q[K]$-resolution $L_\bullet\to M$.
Coinvariants give an exact resolution $(L_\bullet)_N\to M_N$,
free over $\Q[K/N]$, because coinvariants of $\Q[K]$ are
$\Q[K/N]$. Moreover,
\[
 \Q\otimes_{\Q[K]}L_\bullet
 =\Q\otimes_{\Q[K/N]}(L_\bullet)_N.
\]
Computing homology by a resolution of the coefficient module gives
the same answer as resolving the trivial right module: tensor the
two resolutions together, and augment first in either direction.
For each fixed index the other tensor factor is free, so the
augmented rows or columns are exact. Finite elimination along the
diagonals of this first-quadrant double complex identifies its total
homology with both computations. This proves the assertion and its
naturality.
\end{proof}

The simplicial chains carry orientation signs, so the relevant
stabilizers are those of unordered frames with their sign modules.
Unordered resolutions and this use of orientation signs also appear
in the infinite-field setting of Vasilev and Yagunov
\cite{VasilevYagunov2025Unordered}. The following computation supplies
the required statement for the present algebra with residue field
$\F_2$.

\begin{lemma}\label{st:stabilizers}
Let $m\geq13$, put $G=\E_m(R)$, and let $1\leq l\leq4$.
The group $G$ acts transitively on ordered $l$-frames. Put $r=m-l$,
and let $K_l$ be the stabilizer of the unordered frame
$\{e_{r+1},\ldots,e_m\}$. Its complementary-block projection gives
a split exact sequence
\[
 1\longrightarrow N_l\longrightarrow K_l
 \longrightarrow\E_r(R)\longrightarrow1,
\]
where $N_l$ is locally finite. Let $\Q_{\epsilon_l}$ be the
one-dimensional orientation module of this frame. Then
\[
 H_q(K_l;\Q_{\epsilon_l})=
 \begin{cases}
 H_q(\E_{m-1}(R);\Q),&l=1,\\
 0,&2\leq l\leq4,
 \end{cases}
\]
where in the first case the identification is induced by the
complementary inclusion.
\end{lemma}

\begin{proof}
Complete a prescribed ordered frame to a matrix $h\in H_m$, putting
the prescribed columns in the last $l$ positions. By
Lemma~\ref{st:quotient}, the complementary $H_r$ represents every
coset of $\E_m(R)$ in $H_m$, since $r\geq9$. Multiplying $h$ on
the right by an appropriate complementary matrix therefore puts it
in $G$ without changing these columns. This proves transitivity.

In block form the unordered stabilizer in $H_m$ consists of
\[
 \begin{pmatrix}A&0\\B&P\end{pmatrix},
 \qquad A\in H_r,\quad P\text{ a permutation matrix}.
\]
The permutation block is elementary in characteristic two, as is
the strip that removes $B$. The whole matrix belongs to $\E_m(R)$
if and only if its block $A$ does, by injectivity in
Lemma~\ref{st:quotient}. Thus projection onto $A$ has image
$\E_r(R)$ and section $A\mapsto\operatorname{diag}(A,\Id_l)$.
Its kernel is the semidirect product of the additive group
$M_{l\times r}(R)$ with the finite symmetric group $S_l$.

This kernel is locally finite. Given finitely many of its elements,
take all permutation translates of their finitely many strips.
Their span over $\F_2$ is finite, and its semidirect product with
$S_l$ is a finite group containing the given elements.

The orientation character is the sign of $P$. If $l\geq2$, a
transposition in $N_l$ acts as $-1$ on $\Q_{\epsilon_l}$. Thus
its coinvariants are zero: their defining relations include
$-a-a=-2a$. The characteristic of $R$ does not change the rational
orientation sign. Lemma~\ref{st:coinvariants} gives vanishing in
every degree. If $l=1$, the orientation module is trivial and the
same Lemma identifies homology with that of $\E_{m-1}(R)$. The
projection induces this identification, and its complementary
section induces the inverse.
\end{proof}

\begin{proposition}\label{st:h2}
For $n\geq12$, the standard inclusion induces an isomorphism
\[
 H_2(\E_n(R);\Q)\longrightarrow
 H_2(\E_{n+1}(R);\Q).
\]
\end{proposition}

\begin{proof}
Put $m=n+1$, $G=\E_m(R)$, and let $C_\bullet$ be the oriented
simplicial chain complex of $X_m$ over $\Q$, with augmentation
$C_0\to\Q$. Fix a free right resolution $P_\bullet\to\Q$ over
$\Q[G]$. Consider the first-quadrant double complex
\[
 D_{p,q}=P_q\otimes_{\Q[G]}C_p,
\]
with horizontal differential induced by the simplicial boundary,
vertical differential induced by the resolution, and total
differential $d=d_h+(-1)^p d_v$. Write $T_\bullet$ for its total
complex and $A_\bullet=P_\bullet\otimes_{\Q[G]}\Q$.

The augmentation induces a chain map $F:T_\bullet\to A_\bullet$.
We verify explicitly that it induces an isomorphism in degree two.
Adjoin column $p=-1$ by setting $C_{-1}=\Q$. The resulting total
complex has the form
\[
 K_j=T_j\oplus A_{j+1},\qquad
 d_K(t,a)=(d_Tt,F(t)-d_Aa).
\]
Every augmented horizontal row is exact at positions
$p=-1,0,1,2$, by Lemma~\ref{st:frames} and the freeness of each
$P_q$. To kill a cycle of total degree $j\leq2$ in $K_\bullet$,
choose its component of largest resolution degree $q$. That
component is horizontally closed, since no component of larger
$q$ contributes to its horizontal target. Its horizontal position
is $p=j-q\leq2$. Row exactness supplies a horizontal primitive in
position $(p+1,q)$. Subtracting its total boundary removes the
chosen component and introduces only smaller resolution degrees.
Repeating this finite procedure kills the cycle. In particular
$H_1(K_\bullet)=H_2(K_\bullet)=0$. The long exact sequence of
\[
 0\longrightarrow A_\bullet[-1]
 \longrightarrow K_\bullet\longrightarrow T_\bullet
 \longrightarrow0
\]
therefore shows that $F_*:H_2(T_\bullet)\to H_2(A_\bullet)$ is
an isomorphism. Here the connecting map is precisely the
augmentation map $F_*$.

For $0\leq p\leq3$, Lemma~\ref{st:stabilizers} gives a single
orbit of simplices, and the oriented chain module is
\[
 C_p\cong\Q[G]\otimes_{\Q[K_{p+1}]}\Q_{\epsilon_{p+1}}.
\]
The stabilizer is that of the \emph{unordered} simplex, with its
orientation sign on the coefficient line. Tensoring with the right
resolution identifies the vertical column with
$P_\bullet\otimes_{\Q[K_{p+1}]}\Q_{\epsilon_{p+1}}$.
Restriction of $P_\bullet$ to a subgroup is still an exact free
resolution, so the vertical column homology is the stabilizer
homology in Lemma~\ref{st:stabilizers}. Consequently columns one,
two, and three have zero homology in every degree. Column zero
has homology $H_q(\E_{m-1}(R);\Q)$.

Inclusion of column zero induces an isomorphism onto $H_2(T_\bullet)$.
For surjectivity, remove the component in position $(2,0)$ of a
total two-cycle using the vanishing of vertical homology there.
Subtracting the chosen total boundary leaves a cycle whose largest
possible column is one. Remove its component in $(1,1)$ in the
same way. The remaining cycle belongs to column zero.
For injectivity, suppose a column-zero two-cycle bounds a total
three-chain. Remove successively that chain's components in
$(3,0)$, $(2,1)$, and $(1,2)$ by subtracting total boundaries.
At each step the component to remove is vertically closed, because
the boundary is in column zero and all higher components have
already been removed. The requisite vertical homology vanishes.
The resulting bounding chain is in column zero, as required.
Columns of index at least four cannot occur in these cycles or
bounding chains. Equivalently, the only possible incoming
differentials to bidegree $(0,2)$ have sources $(1,2)$, $(2,1)$,
and $(3,0)$, all of which vanish.

For column zero choose the representative vertex $e_m$. Under the
chain identification above, its map through the augmentation is
the map induced by inclusion of its stabilizer into $G$.
The complementary section from $\E_{m-1}(R)$ is exactly
$A\mapsto\operatorname{diag}(A,1)$. Combining the two proved
isomorphisms thus gives the asserted standard stabilization map.
\end{proof}

\begin{proposition}\label{prop:rational-stability}
For $n\geq12$, restriction along $J_n\to J_{n+1}$ induces
an isomorphism, and in particular a surjection,
\[
 \Hom(J_{n+1},\Q)\longrightarrow\Hom(J_n,\Q).
\]
\end{proposition}

\begin{proof}
Proposition~\ref{st:h2} and the bar-cochain identification above
show that restriction gives a bijection on central
extension classes by $\Q$. By Lemma~\ref{st:extensions}, this
is precisely restriction on the displayed Hom groups.
Restriction is a homomorphism of these Hom groups, so its
bijectivity makes it an isomorphism.
\end{proof}
